\documentclass[../../main.tex]{subfiles}
\begin{document}

% 本章导言（PDF p65-66）
微积分的中心思想是线性逼近。这在欧几里得空间微积分的学习中反复出现：例如，一元函数可以用其切线逼近，$\mathbb{R}^{n}$ 中的参数曲线可以用其速度向量逼近，$\mathbb{R}^{3}$ 中的曲面可以用其切平面逼近，或从 $\mathbb{R}^{n}$ 到 $\mathbb{R}^{m}$ 的映射可以用其全导数逼近（见附录~C）。

为使流形上的微积分有意义，我们需要引入流形在一点处的切空间，我们可以把它看作流形在该点附近的某种"线性模型"。由于光滑流形定义的抽象性，这需要一番功夫，我们在本章完成。

我们首先研究具体得多的对象：$\mathbb{R}^{n}$ 中的几何切向量，它们可以看作附着在点上的"箭头"。由于光滑流形的定义是围绕辨认哪些函数光滑这一想法建立的，几何切向量中可推广的性质，是它作为"方向导数"作用于光滑函数的行为。我们在本章第一节证明的关键观察是：取方向导数的过程，在几何切向量与满足乘积法则的从 $C^{\infty}(\mathbb{R}^{n})$ 到 $\mathbb{R}$ 的线性映射之间，给出了一个自然的一一对应。（这类映射称为导子。）以此为动机，我们随后把光滑流形上的切向量定义为 $C^{\infty}(M)$ 在一点处的一个导子。

在本章第二节，我们说明流形之间的一个光滑映射如何在切空间之间诱导一个线性映射，称为该映射的微分，它推广了欧几里得空间之间映射的全导数。这使我们能够把切向量的抽象定义与具体几何图景联系起来：任意光滑坐标图 $(U,\varphi)$ 都给出从 $p$ 处 $M$ 的切空间到 $\varphi(p)$ 处 $\mathbb{R}^{n}$ 的切空间的一个自然同构，而后者又同构于 $\varphi(p)$ 处几何切向量的空间。于是任意光滑坐标图都为每个切空间给出一组基。利用这个同构，我们说明如何在这样一组基中作具体计算。基于这些坐标计算，我们说明一个光滑流形所有点处的所有切空间如何能被"粘合"在一起，形成一个新的流形，称为原流形的切丛。

接着我们说明一条光滑曲线如何在每一点确定一个切向量，称为它的速度，它可以看作 $C^{\infty}(M)$ 中沿曲线取各函数导数的那个导子。

在最后两节中，我们讨论并比较定义切空间的其他几种方法，并简要概述范畴论的术语，它把切空间和微分置于一个更大的背景中。

\section{切向量}\label{section:ch03-切向量}

想象欧几里得空间中的一个流形——例如 $\mathbb{R}^{n}$ 中的单位球面 $S^{n-1}$。我们说 $S^{n-1}$ 一点处的一个"切向量"是什么意思？在回答这个问题之前，我们必须正视我们思考 $\mathbb{R}^{n}$ 元素方式中的一种二分法。一方面，我们通常把它们看作空间中的点，其唯一性质是位置，由坐标 $(x^{1},\dots,x^{n})$ 表达。另一方面，做微积分时我们有时反而把它们看作向量，即具有大小和方向、但其位置无关的对象。向量 $v=v^{i}e_{i}$（其中 $e_{i}$ 表示第 $i$ 个标准基向量）可以看作一个初始点在 $\mathbb{R}^{n}$ 中任意处的箭头；从向量的观点看，唯一相关的只是它指向哪个方向、有多长。

我们这里真正想到的，是在每一点处都有一个单独的 $\mathbb{R}^{n}$ 拷贝。例如，当我们谈论在点 $a$ 处与球面相切的向量时，我们把它们想象为生活在一个原点平移到 $a$ 的 $\mathbb{R}^{n}$ 拷贝中。

\subsection{几何切向量}\label{subsection:ch03-几何切向量}

这是欧几里得空间中切向量的一个初步定义。给定一点 $a\in\mathbb{R}^{n}$，让我们把 $a$ 处 $\mathbb{R}^{n}$ 的几何切空间定义为集合 $\mathbb{R}^{n}_{a}=\{a\}\times\mathbb{R}^{n}=\{(a,v):v\in\mathbb{R}^{n}\}$。$\mathbb{R}^{n}$ 中的一个几何切向量是某个 $a\in\mathbb{R}^{n}$ 处 $\mathbb{R}^{n}_{a}$ 的一个元素。作为记号，我们把 $(a,v)$ 缩写为 $v_{a}$（有时为清晰起见也记作 $v|_{a}$，例如当 $v$ 本身有下标时）。我们把 $v_{a}$ 想象为初始点在 $a$ 的向量 $v$（图~\cref{figure:ch03-3.1}）。集合 $\mathbb{R}^{n}_{a}$ 在如下自然运算下是一个实向量空间：
\begin{align*}
v_{a}+w_{a}=(v+w)_{a},\qquad c(v_{a})=(cv)_{a}.
\end{align*}
向量 $e_{i}|_{a}$（$i=1,\dots,n$）是 $\mathbb{R}^{n}_{a}$ 的一组基。事实上，作为向量空间，$\mathbb{R}^{n}_{a}$ 本质上与 $\mathbb{R}^{n}$ 本身相同；我们加上下标 $a$ 的唯一理由，是使不同点 $a,b$ 处的几何切空间 $\mathbb{R}^{n}_{a}$ 与 $\mathbb{R}^{n}_{b}$ 成为不相交集合。

\begin{figure}[ht]
\centering
\includegraphics[width=0.7\textwidth]{fig-03-1.png}
\caption{几何切空间}\label{figure:ch03-3.1}
\end{figure}

有了这个定义，我们可以把 $S^{n-1}$ 在点 $a\in S^{n-1}$ 处的切空间看作 $\mathbb{R}^{n}_{a}$ 的某个子空间（图~\cref{figure:ch03-3.2}），即与过 $a$ 的径向单位向量正交的向量空间，这里用 $\mathbb{R}^{n}_{a}$ 通过自然同构 $\mathbb{R}^{n}\cong\mathbb{R}^{n}_{a}$ 从 $\mathbb{R}^{n}$ 继承的内积。然而，这个定义的问题在于它没有给我们任何线索，告诉我们如何在任意光滑流形上定义切向量，因为那里没有外围欧几里得空间。所以我们需要寻找切向量的另一种刻画，使之在流形上有意义。

\begin{figure}[ht]
\centering
\includegraphics[width=0.7\textwidth]{fig-03-2.png}
\caption{$S^{n-1}$ 的切空间}\label{figure:ch03-3.2}
\end{figure}

我们在光滑流形上目前拥有的、可用的东西只有光滑函数、光滑映射和光滑坐标图。几何切向量提供的一样东西，是取函数方向导数的手段。例如，任意几何切向量 $v_{a}\in\mathbb{R}^{n}_{a}$ 给出一个映射 $D_{v|_{a}}:C^{\infty}(\mathbb{R}^{n})\to\mathbb{R}$，它取 $a$ 处沿方向 $v$ 的方向导数：
\begin{align}\label{align:ch03-方向导数}
D_{v|_{a}}f=D_{v}f(a)=\left.\frac{\mathrm{d}}{\mathrm{d}t}\right|_{t=0}f(a+tv).
\end{align}
这个运算在 $\mathbb{R}$ 上线性并满足乘积法则：
\begin{align}\label{align:ch03-方向导数乘积法则}
D_{v|_{a}}(fg)=f(a)D_{v|_{a}}g+g(a)D_{v|_{a}}f.
\end{align}
若 $v_{a}=v^{i}e_{i}|_{a}$ 用标准基表出，则由链式法则，$D_{v|_{a}}f$ 可以更具体地写成
\begin{align*}
D_{v|_{a}}f=v^{i}\frac{\partial f}{\partial x^{i}}(a).
\end{align*}
（这里我们照例使用求和约定，所以右端表达式理解为对 $i=1,\dots,n$ 求和。如果我们规定"分母中"的上指标被看作下指标，这个和与我们的指标约定一致。）例如，若 $v_{a}=e_{j}|_{a}$，则
\begin{align*}
D_{v|_{a}}f=\frac{\partial f}{\partial x^{j}}(a).
\end{align*}

带着这一构造，我们给出如下定义。若 $a$ 是 $\mathbb{R}^{n}$ 的一点，称一个映射 $w:C^{\infty}(\mathbb{R}^{n})\to\mathbb{R}$ 为 $a$ 处的一个导子，如果它在 $\mathbb{R}$ 上线性并满足如下乘积法则：
\begin{align}\label{align:ch03-导子乘积法则}
w(fg)=f(a)wg+g(a)wf.
\end{align}
令 $T_{a}\mathbb{R}^{n}$ 表示 $C^{\infty}(\mathbb{R}^{n})$ 在 $a$ 处所有导子所成集合。显然，$T_{a}\mathbb{R}^{n}$ 在运算
\begin{align*}
(w_{1}+w_{2})f=w_{1}f+w_{2}f,\qquad (cw)f=c(wf)
\end{align*}
下是一个向量空间。

关于 $T_{a}\mathbb{R}^{n}$，最重要（也许多少有些令人惊讶）的事实是它有限维，并且实际上自然同构于我们上面定义的几何切空间 $\mathbb{R}^{n}_{a}$。证明基于下面的引理。

\begin{lemma}[导子的性质]\label{lemma:ch03-导子性质}
设 $a\in\mathbb{R}^{n}$，$w\in T_{a}\mathbb{R}^{n}$，$f,g\in C^{\infty}(\mathbb{R}^{n})$。
\begin{enumerate}[(a)]
\item 若 $f$ 是常值函数，则 $wf=0$。
\item 若 $f(a)=g(a)=0$，则 $w(fg)=0$。
\end{enumerate}
\end{lemma}

\begin{proof}
只需对常值函数 $f_{1}(x)\equiv 1$ 证明 (a)，因为 $f(x)\equiv c$ 蕴涵由线性性 $wf=w(cf_{1})=cwf_{1}=0$。对 $f_{1}$，乘积法则给出
\begin{align*}
wf_{1}=w(f_{1}f_{1})=f_{1}(a)wf_{1}+f_{1}(a)wf_{1}=2wf_{1},
\end{align*}
这蕴涵 $wf_{1}=0$。类似地，(b) 也由乘积法则推出：
\begin{align*}
w(fg)=f(a)wg+g(a)wf=0+0=0.
\end{align*}
\end{proof}

下一个命题表明 $a$ 处的导子与几何切向量一一对应。

\begin{proposition}\label{proposition:ch03-导子与几何切向量}
设 $a\in\mathbb{R}^{n}$。
\begin{enumerate}[(a)]
\item 对每个几何切向量 $v_{a}\in\mathbb{R}^{n}_{a}$，由~\eqref{align:ch03-方向导数} 定义的映射 $D_{v|_{a}}:C^{\infty}(\mathbb{R}^{n})\to\mathbb{R}$ 是 $a$ 处的一个导子。
\item 映射 $v_{a}\mapsto D_{v|_{a}}$ 是从 $\mathbb{R}^{n}_{a}$ 到 $T_{a}\mathbb{R}^{n}$ 的一个同构。
\end{enumerate}
\end{proposition}

\begin{proof}
$D_{v|_{a}}$ 是 $a$ 处导子这一事实，是乘积法则~\eqref{align:ch03-方向导数乘积法则} 的直接推论。

为证明映射 $v_{a}\mapsto D_{v|_{a}}$ 是同构，我们首先注意它线性，这容易验证。为看出它单射，假设 $v_{a}\in\mathbb{R}^{n}_{a}$ 具有 $D_{v|_{a}}$ 是零导子的性质。用标准基把 $v_{a}$ 写成 $v_{a}=v^{i}e_{i}|_{a}$，并取 $f$ 为第 $j$ 个坐标函数 $x^{j}:\mathbb{R}^{n}\to\mathbb{R}$（看作 $\mathbb{R}^{n}$ 上的光滑函数），我们得到
\begin{align*}
0=D_{v|_{a}}\bigl(x^{j}\bigr)=v^{i}\frac{\partial}{\partial x^{i}}\bigl(x^{j}\bigr)\bigg|_{x=a}=v^{j},
\end{align*}
其中最后一个等号成立，因为 $\partial x^{j}/\partial x^{i}=0$（除 $i=j$ 时等于 $1$）。由于这对每个 $j$ 都成立，可得 $v_{a}$ 是零向量。

为证明满射，任取 $w\in T_{a}\mathbb{R}^{n}$。受上一段计算启发，我们定义 $v=v^{i}e_{i}$，其中实数 $v^{1},\dots,v^{n}$ 由 $v^{i}=w(x^{i})$ 给出。我们将证明 $w=D_{v|_{a}}$。

为看出这一点，令 $f$ 为 $\mathbb{R}^{n}$ 上任意光滑实值函数。由 Taylor 定理（定理~C.15），我们可以写
\begin{align*}
f(x)=f(a)+\sum_{i=1}^{n}\frac{\partial f}{\partial x^{i}}(a)\bigl(x^{i}-a^{i}\bigr)
+\sum_{i,j=1}^{n}\bigl(x^{i}-a^{i}\bigr)\bigl(x^{j}-a^{j}\bigr)\int_{0}^{1}(1-t)\frac{\partial^{2}f}{\partial x^{i}\partial x^{j}}\bigl(a+t(x-a)\bigr)\,\mathrm{d}t.
\end{align*}
注意上面最后一个和中每一项都是 $x$ 的两个在 $x=a$ 处为零的光滑函数的乘积：一个是 $(x^{i}-a^{i})$，另一个是 $(x^{j}-a^{j})$ 乘以该积分。由引理~\cref{lemma:ch03-导子性质}(b)，导子 $w$ 零化整个这个和。于是
\begin{align*}
wf=w\bigl(f(a)\bigr)+\sum_{i=1}^{n}w\left(\frac{\partial f}{\partial x^{i}}(a)\bigl(x^{i}-a^{i}\bigr)\right)
=0+\sum_{i=1}^{n}\frac{\partial f}{\partial x^{i}}(a)\bigl(w(x^{i})-w(a^{i})\bigr)
=\sum_{i=1}^{n}\frac{\partial f}{\partial x^{i}}(a)v^{i}=D_{v|_{a}}f.
\end{align*}
\end{proof}

\begin{corollary}\label{corollary:ch03-偏导基}
对任意 $a\in\mathbb{R}^{n}$，$n$ 个导子
\begin{align*}
\left.\frac{\partial}{\partial x^{1}}\right|_{a},\dots,\left.\frac{\partial}{\partial x^{n}}\right|_{a}
\end{align*}
（由 $\left.\frac{\partial}{\partial x^{i}}\right|_{a}f=\frac{\partial f}{\partial x^{i}}(a)$ 定义）构成 $T_{a}\mathbb{R}^{n}$ 的一组基，因此 $T_{a}\mathbb{R}^{n}$ 的维数为 $n$。
\end{corollary}

\begin{proof}
应用前一命题，并注意 $\left.\frac{\partial}{\partial x^{i}}\right|_{a}=D_{e_{i}|_{a}}$。
\end{proof}

\subsection{流形上的切向量}\label{subsection:ch03-流形上的切向量}

现在我们可以定义流形和带边流形上的切向量了。两种情形的定义相同。设 $M$ 是一个带边或无边的光滑流形，$p$ 是 $M$ 的一点。称一个线性映射 $v:C^{\infty}(M)\to\mathbb{R}$ 为 $p$ 处的一个导子，如果它满足
\begin{align}\label{align:ch03-流形导子}
v(fg)=f(p)vg+g(p)vf\quad\text{对所有 }f,g\in C^{\infty}(M).
\end{align}
$C^{\infty}(M)$ 在 $p$ 处所有导子所成集合记作 $T_{p}M$，是一个向量空间，称为 $M$ 在 $p$ 处的切空间。$T_{p}M$ 的一个元素称为 $p$ 处的一个切向量。

下面的引理是引理~\cref{lemma:ch03-导子性质} 对流形的类比。

\begin{lemma}[流形上切向量的性质]\label{lemma:ch03-流形切向量性质}
设 $M$ 是一个带边或无边的光滑流形，$p\in M$，$v\in T_{p}M$，$f,g\in C^{\infty}(M)$。
\begin{enumerate}[(a)]
\item 若 $f$ 是常值函数，则 $vf=0$。
\item 若 $f(p)=g(p)=0$，则 $v(fg)=0$。
\end{enumerate}
\end{lemma}

\begin{exercise}\label{exercise:ch03-3.5}
证明引理~\cref{lemma:ch03-流形切向量性质}。
\end{exercise}

带着 $\mathbb{R}^{n}$ 中几何切向量的动机，你应当把 $M$ 的切向量想象为与 $M$ 相切、且基点附着在给定处 $M$ 上的"箭头"。关于切向量定理的证明，当然必须基于以导子表述的抽象定义，但你的直觉应尽可能由几何图景引导。

\end{document}